% ============================================================
%  CV de Fatou FALL - compatible Overleaf (compilateur pdfLaTeX)
% ============================================================
\documentclass[10pt,a4paper]{article}

\usepackage[utf8]{inputenc}
\usepackage[T1]{fontenc}
\usepackage[scaled=0.95]{helvet}
\renewcommand{\familydefault}{\sfdefault}
\usepackage[top=1.5cm,bottom=1.5cm,left=1.8cm,right=1.8cm]{geometry}
\usepackage{xcolor}
\usepackage{titlesec}
\usepackage{enumitem}
\usepackage{tabularx}
\usepackage[hidelinks]{hyperref}

\definecolor{bleu}{HTML}{0369A1}
\definecolor{gris}{HTML}{4B5563}

\pagestyle{empty}
\setlength{\parindent}{0pt}
\setlength{\parskip}{0pt}
\linespread{1.08}

% Titres de section : majuscules, bleu, filet en dessous
\titleformat{\section}{\large\bfseries\color{bleu}}{}{0pt}{\MakeUppercase}[\vspace{2pt}{\color{bleu!40}\titlerule[0.6pt]}]
\titlespacing*{\section}{0pt}{10pt}{6pt}

\setlist[itemize]{leftmargin=1.1em, itemsep=1.5pt, topsep=3pt, parsep=0pt, label=\textbullet}

% \poste{Intitulé}{Organisation, lieu}{Dates}
\newcommand{\poste}[3]{%
  {\bfseries #1}\hfill{\color{gris}\itshape #3}\\
  {\color{bleu}#2}\par\vspace{2pt}}

% \projet{Nom}{Description courte}{Technologies}{Statut}
\newcommand{\projet}[4]{%
  \vspace{5pt}{\bfseries #1}, #2\hfill{\color{gris}\small #4}\\
  {\color{gris}\small #3}\par}

\begin{document}

% ------------------------------------------------------------ En-tête
\begin{center}
  {\fontsize{22}{26}\selectfont\bfseries FATOU FALL}\\[4pt]
  {\large\color{bleu} Développeuse Full Stack : web, mobile et données}\\[6pt]
  {\small ffmaman17@gmail.com \;|\; +221 77 822 84 60 \;|\; Dakar, Sénégal \;|\;
  Portfolio : \href{https://fatou-fall-portfolio.onrender.com}{fatou-fall-portfolio.onrender.com}}
\end{center}
\vspace{-2pt}{\color{bleu!40}\rule{\linewidth}{0.6pt}}

% ------------------------------------------------------------ Profil
\section{Profil}
Développeuse Full Stack en fin de double cursus de Master 2, en Génie Logiciel et en Sciences des Données.
Six mois chez Bakeli sur des produits utilisés en production, et une dizaine d'applications que j'ai
conçues, développées et mises en ligne de bout en bout : un logiciel SaaS pour cabinets d'avocats, une
bibliothèque numérique, une boutique mobile avec paiement Wave et Orange Money. Je ne considère pas un
projet terminé tant qu'il ne tourne pas en ligne (Docker, Render, PostgreSQL, e-mails, stockage), et ma
formation en données me permet aussi de transformer des chiffres en tableaux de bord utiles. Je cherche
une équipe où construire, apprendre vite et prendre des responsabilités.

% ------------------------------------------------------------ Expérience
\section{Expérience professionnelle}
\poste{Développeuse Full Stack (stage)}{Bakeli School of Technology, Dakar}{Mars 2026 - Septembre 2026}
{\small\color{gris} Produits : Easy Desk, Easy Event (web et mobile), EasyCoop, Easy Suite \;|\;
Django, DRF, PostgreSQL, React, Expo, Laravel, Filament, Render}
\begin{itemize}
  \item Conception et développement d'une \textbf{plateforme de support client multi-tenant} avec le chatbot
        \textbf{WAGAN AI} : réponses automatiques aux questions fréquentes et escalade vers un agent humain
        par Discord et e-mail.
  \item Intégration du chatbot dans \textbf{Easy Event}, sur le \textbf{web (React)} et sur
        l'\textbf{application mobile (Expo)}, avec identification de l'utilisateur via Laravel Sanctum et
        traitement des escalades dans Filament.
  \item Déploiement de la même brique dans \textbf{EasyCoop}, en coordination avec les développeurs
        frontend et backend du produit.
  \item Intégration de l'\textbf{authentification unique (SSO) Easy Suite} dans EasyCoop (Laravel, React) :
        une seule connexion pour tous les produits de l'écosystème.
  \item Transformation de l'interface d'administration en \textbf{Easy Desk}, centre de support
        multi-projets, mis en production sur Render (API, base PostgreSQL, interface React).
  \item Rédaction du rapport de stage et de la documentation technique.
\end{itemize}

% ------------------------------------------------------------ Projets
\section{Projets (sélection)}

\projet{Lexora}{logiciel SaaS de gestion de cabinets d'avocats}
{Laravel, React, TypeScript, Tailwind CSS, PostgreSQL, Redis, Docker, Cloudflare R2, Brevo, Render}{En ligne}
\begin{itemize}
  \item Plateforme multi-cabinets : chaque cabinet s'inscrit seul, ses données restent isolées des autres
        (isolation couverte par des tests automatisés).
  \item Dossiers, audiences avec renvois automatiques, agenda, tâches en Kanban, portail client.
  \item Factures générées à partir du temps passé (TVA 18 \%, FCFA), paiements Wave, Orange Money et
        virement, rappel par e-mail la veille de chaque audience.
  \item Intégration continue avec GitHub Actions, Docker multi-services (Nginx, PHP-FPM, Redis).
\end{itemize}

\projet{GlowLuxe By Sofa}{application mobile e-commerce de mode}
{Expo (React Native), Django, PostgreSQL, PayTech, Brevo}{En cours}
\begin{itemize}
  \item Application iPhone, Android et web : catalogue, favoris, panier sans compte fusionné à la connexion.
  \item Paiement Wave, Orange Money et carte, avec acompte selon le montant de la commande.
  \item Attribution automatique des livreurs (disponibilité, charge, distance) et suivi en direct sur carte.
\end{itemize}

\projet{Maktaba CATK}{bibliothèque numérique, réalisée en binôme}
{Django, React, Render, Netlify, Cloudflare R2, Brevo}{En ligne}
\begin{itemize}
  \item API Django sur Render, interface React sur Netlify, déploiement automatique à chaque mise à jour.
\end{itemize}

\projet{KSJI}{site d'un cabinet de services juridiques islamiques (projet client)}
{Django, PostgreSQL (Neon), Cloudflare R2, Brevo, Render}{En ligne}
\begin{itemize}
  \item Recueil du besoin, développement, mise en ligne et livraison au client.
\end{itemize}

\projet{Teranga Voyages}{application pour agence de voyage}{Laravel, React, Docker}{En cours}
\begin{itemize}
  \item Offres, réservations et gestion de l'agence ; conteneurisée pour être déployée et commercialisée.
\end{itemize}

\projet{Al-Barakah}{plateforme de Zakat et de dons}{Laravel, Docker}{En cours}
\begin{itemize}
  \item Calcul de la Zakat, dons et campagnes ; cahier des charges et document technique rédigés en LaTeX.
\end{itemize}

\vspace{4pt}
{\bfseries Aussi :} LocaGest, API de location de véhicules (Spring Boot, MySQL, JWT) ; API de gestion de
stock (Spring Boot) ; application de gestion des emplois du temps (projet de fin de Licence).

% ------------------------------------------------------------ Formation
\section{Formation}
\poste{Master 2 Informatique, Génie Logiciel}{Université Iba Der Thiam de Thiès}{2025 - 2026}
\vspace{3pt}
\poste{Master 2 Sciences des Données et Applications}{Université Iba Der Thiam de Thiès, option Statistique, Économétrie et Modélisation}{2024 - 2025}
\vspace{3pt}
\poste{Licence Informatique, Génie Logiciel}{Université Iba Der Thiam de Thiès}{2020 - 2023}
\vspace{3pt}
\poste{Initiation au développement web avec Django}{She Codes}{Septembre - Octobre 2024}

% ------------------------------------------------------------ Compétences
\section{Compétences techniques}
\begin{tabularx}{\linewidth}{@{}>{\bfseries}p{3.6cm}X@{}}
Back-end           & Django, Django REST Framework, Laravel, Spring Boot, API REST \\[2pt]
Front-end, mobile  & React, TypeScript, Tailwind CSS, Vue.js, React Native (Expo), Flutter \\[2pt]
Bases de données   & PostgreSQL, MySQL, MongoDB, SQL Server, Redis, Firebase \\[2pt]
Sécurité           & JWT, Laravel Sanctum, SSO, CORS, protection anti force brute \\[2pt]
Déploiement        & Docker, GitHub Actions, Render, Netlify, Neon, Cloudflare R2, Brevo \\[2pt]
Données            & Python, R (ggplot2), SQL, Power BI, Stata \\[2pt]
Langages           & Python, PHP, JavaScript, TypeScript, Java, R, SQL \\
\end{tabularx}

\section{Compétences transverses}
\begin{itemize}
  \item Analyse du besoin et rédaction de cahiers des charges et de documentations techniques.
  \item Communication : responsable du pôle Communication et Stratégie Digitale d'une association.
  \item Organisation (daily, reporting), travail en équipe, autonomie et sens du résultat.
\end{itemize}

\section{Langues}
Français : courant \;|\; Anglais : usuel

\end{document}
